\documentclass[11pt]{article}
\usepackage[utf8]{inputenc}
\usepackage[T1]{fontenc}
\usepackage[english]{babel}
\usepackage[a4paper,margin=2.5cm]{geometry}
\usepackage{amsmath,amssymb}
\usepackage{booktabs}
\usepackage{enumitem}
\usepackage{xcolor}
\usepackage{hyperref}

\setlist[itemize]{leftmargin=*, itemsep=0.25em, topsep=0.25em}
\setlist[enumerate]{leftmargin=*, itemsep=0.35em, topsep=0.35em}
\newcommand{\loc}[1]{\path{#1}}
\newcommand{\issue}[1]{\paragraph{#1}}
\emergencystretch=3em

\begin{document}

\begin{center}
  {\Large Manuscript Review}\\[0.5em]
  {\normalsize Review date: 2026-06-08}
\end{center}

\section*{Overall Assessment}

The manuscript has a clear and practically relevant core: a staged deep-reinforcement-learning controller for a hybrid-electric-vehicle surrogate environment, with Phase~1 establishing speed tracking and Phase~2 adding NOx and SOC objectives. The strongest parts are the experimental narrative, the multi-seed reporting, and the explicit discussion of TD3 instability and the WLTC/rule-based comparison.

The draft is not yet submission-ready. The main technical risks are not in the broad research direction, but in implementation-facing ambiguity, inconsistent experimental-budget reporting, unclear emission units, inconsistent SOC criteria, and several conclusions that are stronger than the evidence shown. The manuscript also still contains many visible draft markers: 57 \verb|\note{...}| commands, 86 \verb|\hl{...}| highlights, five headings marked \texttt{[MISSING]}, and one literal ``Section X'' reference. These should be removed or resolved before external review.

I explicitly checked the equations, reward and score definitions, symbol definitions, units, appendix interface tables, sign conventions, dimensional consistency, implementation-facing expressions, and the bibliography. No inverse-trigonometric or coordinate-transform formulas are present, so the branch/quadrant check is not applicable; the analogous implementation-facing sign-convention issue is the EM2 torque sign convention described below.

\section*{Highest-Priority Fixes}

\begin{enumerate}
  \item Clarify the emission units and time integration: decide whether NOx is a mass flow in g/s or a per-step mass in g/step, and make the reward, cumulative totals, and mg/km conversions consistent.
  \item Correct the experimental-design narrative: the methods say roughly 20 HPO trials per algorithm, but the results and discussion report 70 PPO, 45 TD3, and 44 SAC trials.
  \item Define one SOC criterion for ``charge-sustaining'' and apply it consistently: final drift, peak drift, and the Phase~2 Optuna tolerance are currently mixed.
  \item Resolve the observation-normalisation pipeline: the text describes both min--max normalisation to $[0,1]$ and SB3 \texttt{VecNormalize} running-statistic normalisation.
  \item Remove all draft markers, fill the missing Contributions and Conclusion sections, and repair the bibliography errors before compiling the final thesis.
\end{enumerate}

\section*{Technical Findings}

\issue{T1. Emission units are inconsistent between the symbol list, appendix, observation table, reward, and results.}
\begin{itemize}
  \item \textbf{Location:} \loc{main.tex:149}; \loc{appendix.tex:22}; \loc{src/04_implementation.tex:233}; \loc{src/04_implementation.tex:304}; \loc{src/04_implementation.tex:312}; \loc{src/05_results.tex:34}; \loc{src/05_results.tex:336}.
  \item \textbf{Problem:} The same tailpipe NOx quantity is introduced as $NO_{x_{tp}}$ with unit g/s in the symbol list and appendix, but the observation vector lists $\dot m_{NO_x}$ with unit g/step. The reward normalises by $c_{NO_x}=0.4$~g/s, while the results report cumulative NOx in g and intensity in mg/km.
  \item \textbf{Why it matters:} With a 2~Hz sampling rate, treating a g/s signal as g/step changes cumulative emissions by a factor of $\Delta t=0.5$~s. This directly affects the Phase~1 and Phase~2 emission totals, the Euro-limit comparisons, and the reward scale used during training.
  \item \textbf{Suggested fix:} Use one convention throughout. If the surrogate output is a mass flow, write for example
  \[
    E_{NO_x}=\sum_{t=1}^{N}\dot m_{NO_x,t}\,\Delta t,\qquad \Delta t=0.5\ \mathrm{s},
  \]
  and label the observation as g/s. If the implementation stores already integrated per-step mass, relabel it as g/step and change the reward normalisation constant accordingly.
\end{itemize}

\issue{T2. The Phase~1 HPO budget is reported inconsistently and weakens algorithm-level comparisons.}
\begin{itemize}
  \item \textbf{Location:} \loc{src/04_implementation.tex:427}; \loc{src/05_results.tex:46}; \loc{src/05_results.tex:71}; \loc{src/05_results.tex:134}; \loc{src/05_results.tex:192}; \loc{src/06_discussion.tex:70}.
  \item \textbf{Problem:} The implementation chapter states that each algorithm had twenty HPO trials, but the results report 70 PPO trials, 45 TD3 trials, and 44 SAC trials. The discussion explains that old Optuna journal entries were accidentally included.
  \item \textbf{Why it matters:} The selected hyperparameters are not products of equal search budgets. PPO had substantially more search opportunities than SAC or TD3, so statements comparing algorithm quality, sample efficiency, and selected configurations are partly confounded by unequal HPO depth.
  \item \textbf{Suggested fix:} Either rerun the HPO with fresh journals and an equal budget, or state in the Results section before the algorithm comparison that the HPO histories were unequal. Reframe the algorithm rankings as conditional on the available search histories rather than as a controlled comparison.
\end{itemize}

\issue{T3. The observation-normalisation pipeline is ambiguous.}
\begin{itemize}
  \item \textbf{Location:} \loc{src/04_implementation.tex:247}; \loc{src/04_implementation.tex:450}; \loc{src/04_implementation.tex:487}.
  \item \textbf{Problem:} The environment section says every observation component is clipped and min--max normalised to $[0,1]$. Later, the infrastructure section says the environment is wrapped in \texttt{VecNormalize}, which maintains running estimates of observation and reward statistics, and that observations are normalised ``based on the min-max values''.
  \item \textbf{Why it matters:} Min--max scaling and \texttt{VecNormalize} are different transformations. If both are applied, the policy sees normalised min--max features that are then re-centred/re-scaled by running statistics. If only one is applied, the current wording misdocuments the actual input distribution. This affects reproducibility, Phase~2 normaliser reuse, and any attempt to deploy the trained policy.
  \item \textbf{Suggested fix:} Add a short pipeline statement such as: raw physical values $\rightarrow$ clipping $\rightarrow$ min--max scaling $\rightarrow$ optional \texttt{VecNormalize} running mean/variance. If \texttt{VecNormalize} is not applied to observations, remove the conflicting text and specify exactly what is saved in \texttt{vec\_normalize.pkl}.
\end{itemize}

\issue{T4. The backpropagation formula mixes hidden-state and pre-activation gradients.}
\begin{itemize}
  \item \textbf{Location:} \loc{src/02_theory.tex:99}; \loc{src/02_theory.tex:101}; \loc{src/02_theory.tex:108}.
  \item \textbf{Problem:} Equation \texttt{backprop-hidden-gradient} is written as a gradient with respect to $\mathbf{h}^{(l-1)}$, but the right-hand side multiplies by $\sigma'^{(l-1)}(\mathbf{z}^{(l-1)})$, which belongs to the gradient with respect to the previous pre-activation, not the previous hidden activation. Equation \texttt{backprop-weight-gradient} is also too schematic to be implementation-useful.
  \item \textbf{Why it matters:} This is a mathematical-definition section. A reader using the equations to understand or reproduce backpropagation may propagate the wrong quantity or misunderstand the role of $\mathbf{z}^{(l)}$.
  \item \textbf{Suggested fix:} Introduce error signals explicitly, for example
  \[
    \boldsymbol{\delta}^{(l)} = \frac{\partial \mathcal{L}}{\partial \mathbf{z}^{(l)}},\qquad
    \frac{\partial \mathcal{L}}{\partial \mathbf{W}^{(l)}} = \boldsymbol{\delta}^{(l)}\left(\mathbf{h}^{(l-1)}\right)^\top,
  \]
  and
  \[
    \boldsymbol{\delta}^{(l-1)} = \left(\mathbf{W}^{(l)}\right)^\top \boldsymbol{\delta}^{(l)} \odot \sigma'^{(l-1)}\!\left(\mathbf{z}^{(l-1)}\right).
  \]
  If the intent is only conceptual, say so and avoid presenting the schematic chain-rule expression as the operational formula.
\end{itemize}

\issue{T5. The reward-ceiling interpretation is mathematically too strong.}
\begin{itemize}
  \item \textbf{Location:} \loc{src/05_results.tex:9}; \loc{src/05_results.tex:11}; \loc{src/05_results.tex:19}.
  \item \textbf{Problem:} The highlighted claim says that a value close to the 3600 reward ceiling implies speed remains within roughly $\pm5$~km/h for the entire cycle. With $\sigma_v=10$~km/h, a constant 5~km/h error gives a per-step reward of $\exp(-0.5(5/10)^2)\approx0.882$, i.e. about 3177 over 3600 steps, not close to the ceiling. Conversely, an average reward near 95\% of the ceiling corresponds to a constant error near 3.2~km/h, but an average does not prove that every timestep stayed within any bound.
  \item \textbf{Why it matters:} The result interpretation turns a reward average into a uniform tracking guarantee. This overstates what the reward metric can establish.
  \item \textbf{Suggested fix:} Replace the claim with a bounded statement: ``A value close to the ceiling indicates small average tracking errors, but transient excursions must be assessed from RMSE/MAE and the trajectory plots.'' If a maximum-error guarantee is desired, report $\max_t |v_{\mathrm{target},t}-v_{Car,t}|$.
\end{itemize}

\issue{T6. The definition of ``charge-sustaining'' drifts across the manuscript.}
\begin{itemize}
  \item \textbf{Location:} \loc{src/04_implementation.tex:499}; \loc{src/04_implementation.tex:502}; \loc{src/05_results.tex:338}; \loc{src/05_results.tex:339}; \loc{src/05_results.tex:491}; \loc{src/06_discussion.tex:16}.
  \item \textbf{Problem:} The Phase~2 Optuna score constrains peak drift, $\max_t |\Delta\mathrm{SOC}_t|$, with $r_{\mathrm{SOC}}=0.05$. The results and discussion sometimes use final absolute drift, $|\Delta\mathrm{SOC}|$, and sometimes call a seed charge-sustaining if $|\Delta\mathrm{SOC}|<0.10$. The manuscript also states that all Phase~2 seeds are charge-sustaining, while the Phase~2 table reports TD3 peak-drift means of $0.137$ for all seeds and $0.123$ even after removing two divergent seeds.
  \item \textbf{Why it matters:} The central contribution is framed as low NOx under a charge-sustaining constraint. If the criterion changes between optimisation, tables, captions, and discussion, the main claim becomes ambiguous and may be false under the stricter peak-drift criterion.
  \item \textbf{Suggested fix:} Define exactly one primary SOC criterion. For example: ``charge-sustaining means final $|\Delta\mathrm{SOC}|<0.05$'' or ``peak $\max_t|\Delta\mathrm{SOC}_t|<0.05$''. Then update every table and claim. If TD3 fails the criterion, restrict the ``all seeds'' claim to PPO/SAC or remove it.
\end{itemize}

\issue{T7. Several notation variants describe the same implementation variables.}
\begin{itemize}
  \item \textbf{Location:} \loc{main.tex:146}; \loc{main.tex:148}; \loc{main.tex:149}; \loc{src/03_methodology.tex:54}; \loc{src/04_implementation.tex:229}; \loc{src/04_implementation.tex:230}; \loc{src/04_implementation.tex:304}; \loc{appendix.tex:49}.
  \item \textbf{Problem:} The target speed appears as $v_{\mathrm{target}}$, $v_{\mathrm{target},t}$, $v_{\text{ref},t}$, and sometimes just $v_t$ for actual speed. Battery state of charge appears as $SoC$, $\mathrm{SOC}$, $\text{SOC}$, and the acronym SOC. Tailpipe NOx appears as $NO_{x_{tp}}$, $\dot m_{NO_x}$, $\tilde E_{NO_x}$, and cumulative $E_{NO_x}$.
  \item \textbf{Why it matters:} These are not just stylistic differences: they connect the symbol list, plant-model outputs, reward code, and reported metrics. Notation drift makes implementation reproduction error-prone.
  \item \textbf{Suggested fix:} Add a notation table for the RL environment and use it consistently. For example, reserve $v_{\mathrm{ref},t}$ for the target trajectory, $v_{Car,t}$ for measured vehicle speed, $\mathrm{SOC}_t$ for battery state, $\dot m_{NO_x,t}$ for mass flow, and $E_{NO_x}$ for integrated mass.
\end{itemize}

\issue{T8. The EM2 torque sign convention is asserted but not defined.}
\begin{itemize}
  \item \textbf{Location:} \loc{src/04_implementation.tex:238}; \loc{src/04_implementation.tex:363}; \loc{src/05_results.tex:168}; \loc{src/05_results.tex:224}.
  \item \textbf{Problem:} The safety filter states that charging corresponds to $\tau_{EM2}<0$ and discharging to $\tau_{EM2}>0$. Later action assessments interpret positive/negative EM2 torque, but no sign convention is defined in the symbol list or action table.
  \item \textbf{Why it matters:} This is a sign convention with direct consequences for the safety filter. Reversing it would make the boundary-protection logic harmful rather than protective.
  \item \textbf{Suggested fix:} Add a sentence and, ideally, a small equation defining the sign convention, e.g. ``positive $\tau_{EM2}$ denotes traction/discharge at the battery and negative $\tau_{EM2}$ denotes regenerative or charging operation,'' if that is indeed the implemented convention.
\end{itemize}

\issue{T9. The NOx-mechanism interpretation overstates monotonicity and causality.}
\begin{itemize}
  \item \textbf{Location:} \loc{src/05_results.tex:507}; \loc{src/05_results.tex:508}; \loc{src/05_results.tex:510}; \loc{src/06_discussion.tex:20}; \loc{src/06_discussion.tex:21}; \loc{src/06_discussion.tex:23}.
  \item \textbf{Problem:} The text says NOx rate rises monotonically with engine torque and then uses this to explain the Phase~2 operating-point shift. The evidence named in the text is a Pearson correlation of 0.84, which supports a positive association but does not establish monotonicity. The sentence beginning ``Assuming that overcharge required additional high-load engine operation...'' is also a fragment and marks a causal assumption rather than a demonstrated mechanism.
  \item \textbf{Why it matters:} The mechanistic explanation is one of the strongest interpretive claims in the Results and Discussion. It should distinguish observed association, plausible mechanism, and proven causal effect.
  \item \textbf{Suggested fix:} Rephrase to ``in the pooled bench/map data, NOx rate is strongly positively associated with torque'' unless a monotonic regression or binned monotonic trend is shown. Replace the fragment with a qualified causal statement such as: ``This suggests that Phase~1 overcharging required additional high-load engine operation, inflating both fuel and NOx.''
\end{itemize}

\issue{T10. The WLTC/rule-based comparison needs a fuller protocol and a fairness caveat.}
\begin{itemize}
  \item \textbf{Location:} \loc{src/05_results.tex:546}; \loc{src/05_results.tex:551}; \loc{src/05_results.tex:581}; \loc{src/05_results.tex:585}; \loc{src/05_results.tex:605}; \loc{src/05_results.tex:631}; \loc{src/06_discussion.tex:45}; \loc{src/06_discussion.tex:46}.
  \item \textbf{Problem:} The WLTC zero-shot and rule-based comparison appears in Results, but the Methodology chapter still contains a note saying to add this comparison. The RL agents start and return to $\mathrm{SOC}\approx0.70$, while the rule-based isoSOC run uses $\mathrm{SOC}\approx0.56$.
  \item \textbf{Why it matters:} If the plant model, battery limits, or controller behaviour depend on absolute SOC, then charge neutrality alone does not guarantee a fair comparison. The best-seed selection criterion is also introduced only locally in Results.
  \item \textbf{Suggested fix:} Add a WLTC evaluation subsection to Methodology. State the initial SOCs, why different absolute SOC levels are comparable, the exact rule-based baseline source, the seed-selection score, and whether the same normalisation/scalers and safety filters are used.
\end{itemize}

\issue{T11. One Phase~1 quantitative statement contradicts the table.}
\begin{itemize}
  \item \textbf{Location:} \loc{src/05_results.tex:32}; \loc{src/05_results.tex:214}.
  \item \textbf{Problem:} The text says SAC has both the lowest mean RMSE and the lowest per-seed variance in Phase~1. The table gives SAC $3.45\pm0.43$~km/h and TD3 $3.81\pm0.29$~km/h, so TD3 has the lower RMSE standard deviation.
  \item \textbf{Why it matters:} This affects the claim that SAC is the most reliable speed tracker in Phase~1.
  \item \textbf{Suggested fix:} Write that SAC has the best mean RMSE, while TD3 has the tightest RMSE dispersion but worse energy/emission dispersion.
\end{itemize}

\issue{T12. The Optuna-vs-grid-search conclusion is stronger than the manuscript supports.}
\begin{itemize}
  \item \textbf{Location:} \loc{src/03_methodology.tex:226}; \loc{src/04_implementation.tex:491}; \loc{src/06_discussion.tex:40}; \loc{src/06_discussion.tex:43}.
  \item \textbf{Problem:} The discussion states that continuous Bayesian optimisation materially outperforms grid search for narrow stable regions. The manuscript describes both procedures, but does not present a controlled grid-versus-Optuna comparison with matched budgets and outcomes.
  \item \textbf{Why it matters:} This is an unsupported methodological conclusion.
  \item \textbf{Suggested fix:} Reframe as: ``Optuna was more targeted than the initial grid and found good configurations within the tested budget.'' Only claim superiority if the grid and Optuna results are tabulated under comparable compute.
\end{itemize}

\section*{Notation and Terminology Drift Check}

\begin{itemize}
  \item \textbf{$v_{\mathrm{target}}$ / $v_{\mathrm{ref}}$ / $v_t$:} First defined in the symbol list as target/reference vehicle velocity. Later the RMSE uses $v_{\mathrm{target},t}$, while the reward defines $\Delta v_t=v_{\text{ref},t}-v_t$. This is a variant rather than a contradiction, but it is implementation-facing and should be standardised.
  \item \textbf{$SoC$ / $\mathrm{SOC}$ / $\text{SOC}$:} First defined as battery state of charge. Later the manuscript alternates typography and alternates between signed final drift, absolute final drift, and peak drift. This is a real conceptual ambiguity because the central ``charge-sustaining'' claim depends on it.
  \item \textbf{$NO_{x_{tp}}$ / $\dot m_{NO_x}$ / $E_{NO_x}$:} First defined as tailpipe NOx mass flow in g/s. Later it is listed as g/step and integrated into cumulative grams and mg/km. This is a unit and implementation ambiguity, not merely a notation preference.
  \item \textbf{$\tau_{EM2}$:} First defined as EM2 torque. Later its sign is used to decide charging/discharging. The sign convention is not verbally attached at first definition and should be added.
\end{itemize}

\section*{Meaningful Editorial and Completeness Findings}

\issue{E1. Visible draft markers and missing sections remain throughout the compiled manuscript.}
\begin{itemize}
  \item \textbf{Location:} Examples include \loc{src/01_introduction.tex:72}, \loc{src/06_discussion.tex:6}, \loc{src/07_conclusion.tex:1}, \loc{src/07_conclusion.tex:4}, \loc{src/07_conclusion.tex:6}, and many highlighted notes across the source.
  \item \textbf{Problem:} The manuscript still contains visible \verb|\note{...}| and \verb|\hl{...}| content, sections marked \texttt{[MISSING]}, and an empty Conclusion chapter.
  \item \textbf{Why it matters:} This is a major professionalism and completeness issue. It also obscures which claims are final and which are internal TODOs.
  \item \textbf{Suggested fix:} Resolve all highlighted notes, remove the \verb|\note| macro from final prose, fill the Contributions and Conclusion sections, and recompile with a search for \verb|\hl{|, \verb|\note{|, \texttt{[MISSING]}, and \texttt{Section X} returning zero hits.
\end{itemize}

\issue{E2. A broken placeholder reference remains in the SAC results.}
\begin{itemize}
  \item \textbf{Location:} \loc{src/05_results.tex:224}.
  \item \textbf{Problem:} The text says the shielding mechanism is explained in ``Section X''.
  \item \textbf{Why it matters:} This is a visible broken cross-reference in a technical interpretation paragraph.
  \item \textbf{Suggested fix:} Replace it with the actual reference to the safety-mechanism paragraph, e.g. \verb|\autoref{sec:impl_environment}| or a more specific label added to the engine-switch cooldown paragraph.
\end{itemize}

\issue{E3. The discussion confuses CO, CO$_2$, fuel, and the single-pollutant limitation.}
\begin{itemize}
  \item \textbf{Location:} \loc{src/06_discussion.tex:48}; \loc{src/06_discussion.tex:50}; \loc{src/06_discussion.tex:88}.
  \item \textbf{Problem:} The text says the engine-always-on and elevated fuel use are consequences of not minimising ``CO emissions'' while the following paragraph discusses fuel and CO$_2$. CO and CO$_2$ are different pollutants/quantities, and neither is identical to fuel consumption.
  \item \textbf{Why it matters:} This blurs the environmental interpretation of the reward design.
  \item \textbf{Suggested fix:} Use ``fuel/CO$_2$ blind spot'' when discussing fuel economy and greenhouse-gas implications. Discuss CO separately only if the model output $CO_{tp}$ is evaluated.
\end{itemize}

\issue{E4. Some claims need citation or removal before final submission.}
\begin{itemize}
  \item \textbf{Location:} Examples include \loc{src/01_introduction.tex:4}, \loc{src/01_introduction.tex:6}, \loc{src/01_introduction.tex:39}, \loc{src/01_introduction.tex:41}, \loc{src/04_implementation.tex:162}, \loc{src/06_discussion.tex:98}, and \loc{src/06_discussion.tex:153}.
  \item \textbf{Problem:} Several factual claims are still highlighted with instructions to cite or check them, including transport-sector emissions, combustion-vehicle percentages, CORNET/FVV context, Euro regulation framing, CTTC/JFF details, carbon-intensity values, UPC codes, and SDG statements.
  \item \textbf{Why it matters:} These are not stylistic placeholders; they are factual support gaps.
  \item \textbf{Suggested fix:} Either provide authoritative citations or remove the claims. Keep unsupported background claims out of the final argument if they are not essential.
\end{itemize}

\section*{Proofing-Sweep Items}

The bundled mechanical proofing scan was run on the compiled PDF and on a combined source-text file. The PDF hits were false positives caused by table-of-contents dot leaders. The source-text hits were mostly comments and legitimate ADS Bibcode strings containing ``..''. No arctangent/branch-convention hit was present. Manual spot-checking still identified the following high-confidence proofing and reference-hygiene issues.

\begin{itemize}
  \item \textbf{Bibliography duplicate keys:} \loc{master_thesis.bib:518} repeats \texttt{noauthor\_soft\_nodate}, and \loc{master_thesis.bib:661} repeats \texttt{wang\_energy\_2025}. \loc{main.blg:6} and \loc{main.blg:10} report repeated-entry errors. Rename the duplicate keys and update citations if needed.
  \item \textbf{Malformed edition output:} \loc{main.bbl:147} prints ``second edition edition'' for Sutton and Barto. Change \loc{master_thesis.bib:352} from \texttt{edition = \{Second edition\}} to a style-compatible value such as \texttt{edition = \{2\}} or \texttt{edition = \{Second\}}.
  \item \textbf{No-author BibTeX warnings:} \loc{main.blg:14}--\loc{main.blg:18} warn that several \texttt{noauthor\_*} entries need an author or key for sorting. Add \texttt{author}, \texttt{organization}, or \texttt{key} fields for the documentation/software references.
  \item \textbf{Typographical cleanup:} High-confidence spelling defects include ``unfeasable'' (\loc{src/01_introduction.tex:24}), ``utlises'' and ``utlise'' (\loc{src/02_theory.tex:7}, \loc{src/02_theory.tex:14}), ``looses'' for ``loses'' (\loc{src/02_theory.tex:14}), ``adressing'' (\loc{src/02_theory.tex:118}), ``directiosn'' (\loc{src/02_theory.tex:255}), ``dependy'' and ``faciliate'' (\loc{src/04_implementation.tex:139}), and ``plateu'' (\loc{src/05_results.tex:156}).
  \item \textbf{Duplicate centering command:} \loc{src/05_results.tex:609}--\loc{src/05_results.tex:611} contains two consecutive \verb|\centering| commands in the same figure. Remove one.
\end{itemize}

\section*{Items Needing External Verification}

These items may be correct, but the manuscript does not yet provide enough support to treat them as established facts.

\begin{itemize}
  \item \textbf{Regulatory comparisons:} The Euro~6 NOx limit, Euro~7 framing, and the interpretation of staircase-cycle mg/km values as regulatory context need authoritative citations and a clear caveat that the staircase cycle is not a certification WLTC/RDE procedure.
  \item \textbf{Current industry practice:} Claims that the automotive industry mainly relies on rule-based EMSs, and that MPC/ECMS generally struggle in the stated ways, should be supported by current sources or narrowed to the cited literature.
  \item \textbf{Carbon-intensity values:} The Spain and Germany electricity-carbon-intensity values in the sustainability analysis are time-dependent snapshot values; cite the source, date, and method.
  \item \textbf{Project and institutional context:} CORNET, FVV, UPC/IDEAI, CTTC/JFF, and UPC ethics/SDG statements need exact institutional references if retained.
  \item \textbf{Rule-based baseline:} The WLTC rule-based benchmark from Frey et al. should be tied to a precise source, scenario, initial SOC, and measurement/calculation protocol.
\end{itemize}

\end{document}